\documentclass[../../../include/open-logic-section]{subfiles}

\begin{document}

\olfileid[id]{sth}{story}{urelements}
\olsection{Urelemen atau Bukan?}

Dalam beberapa bab berikutnya, kita akan mencoba menurunkan aksioma-aksioma
dari konsepsi kumulatif-iteratif tentang himpunan. Namun, sebelum melangkah
lebih jauh, kita perlu mengatakan sesuatu lagi tentang \emph{urelemen}.

Gambaran dalam \olref[approach]{sec} hanya memperbolehkan kita membentuk
himpunan baru dari \emph{himpunan} lama. Akan tetapi, kita mungkin ingin
memperbolehkan \emph{bukan-himpunan} tertentu---sapi, babi, butir pasir, atau
apa pun---menjadi !!{element}s suatu himpunan. Dalam hal itu, kita akan memulai
dengan unsur dasar tertentu, yaitu \emph{urelemen}, lalu mengatakan bahwa pada
setiap tahap $S$ kita membentuk ``semua himpunan yang mungkin'' yang terdiri
atas urelemen atau himpunan yang dibentuk pada tahap sebelum $S$ (dalam
kombinasi apa pun). Gambaran yang dihasilkan akan lebih menyerupai ini:
\begin{center}
	\begin{tikzpicture}[scale=0.6]
	\tikzset{cut_here/.style={densely dotted}}
	\draw (-4,6) -- (-1, 0) -- (1,0) -- (4,6);
	\draw[cut_here] (-5,8)--(-4,6);
	\draw[cut_here] (5,8)--(4,6);
	\draw(-1.5, 1)--(1.5, 1);
	\draw(-2, 2)--(2, 2);
	\draw(-2.5, 3)--(2.5,3);
	\draw(-3, 4)--(3, 4);
	\draw(-3.5, 5)--(3.5, 5);
	\draw(-4, 6)--(4, 6);
	\node[label] at (2, 0) {\small 0};
	\node[label] at (2.5, 1) {\small 1};
	\node[label] at (3, 2) {\small 2};
	\node[label] at (3.5, 3) {\small 3};
	\node[label] at (4, 4) {\small 4};
	\node[label] at (4.5, 5) {\small 5};
	\node[label] at (5, 6) {\small 6};
	\end{tikzpicture}
\end{center}
Jadi, sekarang ada keputusan yang harus kita ambil: \emph{Haruskah kita
memperbolehkan urelemen?}

Secara filosofis, memasukkan urelemen ke dalam teori kita masuk akal. Alasan
utamanya ialah agar teori himpunan kita \emph{dapat diterapkan}. Sebagai
ilustrasi, ingat dari \olref[sfr][siz][]{chap} bahwa kita mengatakan dua
himpunan $A$ dan~$B$ berukuran sama, yakni $\cardeq{A}{B}$, jika dan hanya
jika terdapat suatu bijeksi di antara keduanya. Sekarang, jika sapi di padang dan babi di kandang
masing-masing membentuk himpunan, kita dapat memberikan perlakuan
teoretis-himpunan terhadap klaim ``jumlah sapi sama banyaknya dengan jumlah
babi''. Namun, jika kita melarang urelemen sehingga sapi dan babi \emph{tidak}
membentuk himpunan, perlakuan teoretis-himpunan tersebut tidak akan tersedia.
Memang, kita tidak akan mempunyai cara langsung untuk menerapkan teori
himpunan pada apa pun selain himpunan itu sendiri. (Untuk alasan lain yang
mendukung urelemen, lihat \citealt[pp.~vi, 24, 50--1]{Potter2004}.)

Namun, secara matematis, urelemen cukup jarang diperbolehkan. Sebagian
alasannya ialah bahwa teori himpunan \emph{sedikit sekali} lebih mudah
dirumuskan tanpa urelemen. Akan tetapi, sesekali kita menemukan pembenaran
yang lebih menarik untuk mengeluarkan urelemen dari teori himpunan:
\begin{quote}
	Sesuai dengan keyakinan bahwa teori himpunan merupakan fondasi matematika,
	kita seharusnya dapat mencakup seluruh matematika hanya dengan membicarakan
	himpunan; karena itu, variabel kita tidak boleh menjangkau objek
	seperti sapi dan babi.
	%But if $C$ is a cow, $\{C\}$ is a set, but not a legitimate mathematical object.
	\citep[p.~8]{Kunen1980}
\end{quote}
Jadi, fokus pada keterterapan menyarankan agar urelemen \emph{dimasukkan};
fokus pada tujuan fondasional yang reduktif (mereduksi matematika menjadi
teori himpunan murni) mungkin menyarankan agar urelemen \emph{dikeluarkan}.
Sedikit kemalasan juga mengarah pada pengeluaran urelemen.

Kita akan mengikuti jalan yang paling malas. Namun, sebagian pilihan ini juga
memiliki pembenaran pedagogis. Tujuan kita ialah memperkenalkan kepada Anda
unsur-unsur teori himpunan yang diperlukan untuk mulai mempelajari filsafat
teori himpunan. Dan sebagian besar literatur filosofis tersebut membahas
teori himpunan yang dirumuskan \emph{tanpa} urelemen. Jadi, buku ini barangkali
akan lebih berguna jika mengikuti literatur tersebut dengan cukup dekat.

\end{document}
